\section{A geometric family with exponentially many lower forests}
\label{sec:forest-separation-family}

The small closed-walk example extends to a family in which the candidate
forest space is exponential, although delayed discovery is linear in the
input size.  This does not establish a lower bound on all competing
algorithms: the particular family also admits a direct block decomposition.
Its purpose is to prove a strict, scalable difference between a canonical
forest policy, unpruned forest enumeration, and the new complete strategy.

\begin{theorem}[Theta-block separation family]
\label{thm:theta-family}
For every $h\ge1$, there is a connected graph $G_h$ with $h+1$ vertices and
$3h$ edges, and $2h$ specified closed attaching walks, with the following
properties.
\begin{enumerate}
\item $G_h$ has exactly $3^h$ spanning trees.  Exactly $2^h$ of them permit
complete raw positive-singleton cancellation of all $2h$ rows.
\item The source-order canonical forest leaves every row in a stopping
set.  Full-matrix peeling deletes all rows, using $5h$ nonzero incidences,
and then obtains a disjoint spanning forest in $O(h)$ indexed work.
\item These unsigned attaching incidences admit a genuine geometric
handle realization of a three-ball.  This existence statement does not
assert that the repository's normal-cut front end produces that handle
structure.
\end{enumerate}
\end{theorem}
\begin{proof}
Take vertices $v_0,v_1,\ldots,v_h$.  For each $i$, join $v_0$ to $v_i$ by
three parallel edges $a_i,b_i,c_i$, in this source order.  The two closed
walks in block $i$ have words
\[
 a_ic_i^{-1}b_ic_i^{-1},\qquad b_ic_i^{-1}.
\]
Their raw incidence block, with columns $a_i,b_i,c_i$, is
\[
 N_i=\begin{pmatrix}1&1&2\\0&1&1\end{pmatrix}.
\]
There are no incidences between blocks.  Every spanning tree must contain
exactly one edge from each triple, so there are $3^h$ trees.

Choosing $a_i$ in a block leaves
$\left(\begin{smallmatrix}1&2\\1&1\end{smallmatrix}\right)$, a nonempty
stopping set.  No operations in other blocks can change that obstruction.
Choosing $b_i$ or $c_i$ leaves respectively
\[
 \begin{pmatrix}1&2\\0&1\end{pmatrix},\qquad
 \begin{pmatrix}1&1\\0&1\end{pmatrix},
\]
each completely column-peelable.  Thus success is equivalent to choosing
one of $b_i,c_i$ independently for every block, giving exactly $2^h$
successful trees.  The canonical tree chooses every $a_i$, so its entire
$2h$-row matrix is blocked.  In the full matrix, column $a_i$ first removes
the first row of block $i$, after which $b_i$ or $c_i$ removes the second.
There are $6h$ literal occurrences and $5h$ nonzero incidences, and each is
processed only a bounded number of times.

For geometric realization, begin with two zero-handles and three
one-handles $a_i,b_i,c_i$ joining them.  A zero-handle boundary with its
three one-handle feet removed is a pair of pants.  Choose two disjoint
seams joining the $a_i$ foot to the $c_i$ foot and the $b_i$ foot to the
$c_i$ foot.  Choose compatible seams at the other zero-handle.  Distinct
product strips through the one-handles join the seams into two disjoint
simple attaching curves $x_i,y_i$ whose literal passages are
$a_ic_i^{-1}$ and $b_ic_i^{-1}$.  Contracting $c_i$ gives a single
zero-handle, and the two curves cancel the remaining one-handles, so
attaching two-handles along this system gives a ball.

Slide the two-handle attached along $x_i$ over the one attached along
$y_i$.  Realize the slide by an embedded band in the boundary of the
central zero-handle, with interior disjoint from both curves and every
one-handle foot, and use a parallel copy of $y_i$.  Such a band exists in
the complement of the two chosen seams.  The slid curve and the original
$y_i$ remain disjoint and simple.  The based word can depend on orientation and basepoint conventions;
the theorem asserts realization of the unsigned incidence matrix, for
which either orientation of the parallel copy has the same count.  Since
the band itself stays in the zero-handle, its raw passages are precisely
the sum of the passages of $x_i$ and $y_i$, namely $(1,1,2)$.  The second
row remains $(0,1,1)$.  A geometric two-handle slide preserves the
three-manifold, so this is still a handle decomposition of a ball.

Place $h$ copies in disjoint regions of one central zero-handle, with
separate outer zero-handles.  Equivalently, take boundary connected sums
of the constructed balls and merge their central zero-handles.  The
result is a three-ball with exactly the graph $G_h$ and the stated
unsigned block incidence matrix.  The proof is an explicit handle
construction; converting it to the particular authenticated
normal-parallelity source schema remains separate implementation work.
\end{proof}

\begin{proposition}[Cost of a specified unpruned enumeration]
For the family above, enumerate spanning trees in lexicographic order
$(s_1,\ldots,s_h)\in\{0,1,2\}^h$, with choices $a_i,b_i,c_i$ respectively,
and run the column-leaf test separately on each tree.  The first success
occurs after
\[
 \frac{3^h+1}{2}
\]
tested trees.
\end{proposition}
\begin{proof}
The lexicographically first successful vector is $(1,\ldots,1)$.  Its
zero-based ternary index is
$\sum_{j=0}^{h-1}3^j=(3^h-1)/2$.  Including this final successful trial
gives the displayed count.
\end{proof}

A block-aware solver does not need this enumeration, and neither does
the deferred-forest theorem.  The exact count concerns the stated
baseline, not an intrinsic exponential lower bound.  The saved benchmark
driver reports all raw timings for this baseline, for deferred discovery
on the same family, and for both native core policies on the common
success subset.  Source geometry preparation, dual reduction, and
certificate replay are excluded from those kernel timings and must not
be conflated with end-to-end recognition time.
